\begin{textsample}{POS Dim 6 – vem\_ed\_32 – Score 4.00 – vem\_ed\_32/t170.txt}  \label{ex:f6_pos_006}
Mesa-redonda virtual debate \textbf{benefícios} dos projetos COIL Maria Claudia Nunes Delfino e Neusa Haruka Sezaki Gritti, coordenadoras de Projetos Colaborativos Internacionais (PCIs) em inglês na Área de Apoio à Internacionalização do Ensino Superior da Divisão de Extensão e Pesquisa (Depes) da Coordenadoria de Ensino Superior de Graduação (CGESG) do Centro Paula Souza (CPS), e Tálita Guarino, professora de inglês da Fatec Guaratinguetá, participaram de mesa-redonda virtual no XIV Congresso Internacional de Ciência, Tecnologia e Desenvolvimento (CICTED), evento gratuito promovido pela Universidade de Taubaté (Unitau) em outubro de 2025.
A mesa-redonda foi composta ainda por Anoush Ayunts, da Yerevan State University (Armênia) e contou com a moderação de Adriana Leônidas de Oliveira (Unitau).
Neusa trouxe as \textbf{definições} de Intercâmbio Virtual e da abordagem COIL.
Claudia resumiu as fases dos projetos COIL e em sua articulação para a busca de solu- ções para problemas locais ou globais. “COIL se adapta a múltiplas disciplinas: desde administração até estudos ambientais e ciências da saúde, mas o que une todas essas experiências é o foco na colaboração e compreensão intercultural”.
As \textbf{palestrantes} apresentaram os \textbf{benefícios} culturais e linguísticos, para estu- dantes de graduação, dos projetos COIL – que, nas Fatecs do CPS, recebem o nome de PCIs.
A professora Anoush realiza, desde 2021, um PCI entre a Yerevan State University e Fatec São José do Rio Preto, com a parceria da professora Edilene Gasparini Fernandes, que incluiu Webquests no projeto (leia Artigo de Opinião em VEm 25: https://tinyurl.com/5d9pfp6k).
De 2022 em diante, o projeto passou a abordar Objetivos de Desenvolvimento Sustentável da ONU e se es- tendeu para as \textbf{unidades} de Bragança Paulista, Garça, Jales, Mogi das Cruzes, Pompeia e Guaratinguetá (com Tálita Guarino).
Anoush comentou sobre uma publicação derivada dessa experiência, na Revista CBTecLE, o depoimento do aluno Israel Nascimento para o \textbf{canal} da UNICollaboration no YouTube e a publicação no site do CPS sobre o PCI.
Tálita Guarino compartilhou depoimentos em vídeo de três estudantes e elencou \textbf{benefícios} do PCI para alunos, professores e instituições de ensino superior (leia mais sobre o assunto no Artigo de Opinião).
A professora identificou quatro be- nefícios principais a partir dos relatos dos discentes: Colaboração intercultural – aprender a respeitar e valorizar diferentes perspectivas; Desenvolvimento linguístico - usar o inglês como ferramenta para conexão, não para perfeição; Competência digital – dominar ferramentas tecnológicas; Aprendizagem socioemocional – empatia, tolerância, flexibilidade e confiança. “Os estudantes mudaram seu mindset – pararam de se preocupar com erros e começaram a focar em interação significativa”, comentou Tálita.
Ela elencou ainda os \textbf{benefícios} para os professores, em uma transformação pedagógica: de “entregadores de conteúdo” para facilitadores de aprendizagem.
Isso envolve um planejamento intencional junto com o parceiro da IES internacional, gerenciamento das diferenças culturais e de fusos horários.
Para as instituições, o COIL promove Internacionalização em Casa de forma \textbf{autêntica}.

% matched lemmas: autêntico, benefício, canal, definição, palestrante, unidade
\end{textsample}
